% Notes of the Montreal seminar, 1970 (cote 129, typed part): archive cote 129, site page 31 (PDF page 32).
% typescript. Draft compared: none.
% AI-transcribed by Claude Opus 5.5 (claude-opus-5-5, adaptive thinking, effort high), Message Batches API, batch pilotnext (msgbatch_01A9sdu8SYq23bEtqX9tUncQ), 2026-11-13; not yet checked by eye.
\scanpage{129/31}
\mspage{\hand{II.3}}

\[ F : \mathbb{G}_{a|S} \to \mathbb{G}_{a|S} \qquad x \to x^{p} . \]

On désigne par $\alpha_{p|S}$ le noyau de $F$ . C'est un groupe libre fini de rang
$p$ sur $S$ , d'algèbre affine isomorphe à $\underline{\mathcal{O}}_{S}[T]/T^{p}$ .

1.4. \uline{Dualité de Cartier}

Pour tout schéma en groupes (commutatif) $G$ sur $S$ , on pose :
\[ G^{*} = \underline{\mathrm{Hom}}_{S\text{-}gr}(G,\mathbb{G}_{m|S}) . \]

\margin{$x \mapsto$} On définit ainsi un foncteur contravariant $G \over{\to}{\mapsto} G^{*}$ de la catégorie des
groupes sur $S$ dans elle-même ; ce foncteur commute aux changements de base.

Si $G = \mathrm{Spec}(A)$ est fini localement libre sur $S$ , on montre
(SGA 3, VII$_{A}$, 3.3.1) que l'on a
\[ G^{*} = \mathrm{Spec}(\overset{\hand{\vee}}{A}) \]
avec
\[ \overset{\hand{\vee}}{A} = \underline{\mathrm{Hom}}_{\underline{\mathcal{O}}_{S}\text{-mod}}(A,\underline{\mathcal{O}}_{S}) , \]
\margin{$\times$} la structure de \deleted{bi-algèbre}\inserted{bigèbre} sur $\overset{\hand{\vee}}{A}$ étant définie à l'aide des transposés des
morphismes $\Delta$ et $\mu$ de 1.1. Alors, $G^{*}$ est donc aussi un groupe fini lo-%
calement libre sur $S$ , et $G^{*}$ est de même rang que $G$ . De plus, l'homo-%
\margin{$\times$} morphisme fonctoriel canonique $G \to (G^{*})^{\hand{*}}$ déduit de la définition de $G^{*}$
est défini par l'isomorphisme $A \simeq (\overset{\hand{\vee}}{A})^{\vee}$ , et est donc un isomorphisme :
\[ G \overset{\sim}{\to} (G^{*})^{*} . \]

On définit donc ainsi une anti-équivalence de la catégorie des
$S$-schémas en groupes (commutatifs) finis localement libres avec elle-même,
appelée \uline{dualité de Cartier}.

\mspage{\hand{31}}
